\chapter{The Quadratic Form as an Operator}
\label{ch:operator}

The point that does the most work in practice is the easiest to miss, because every
textbook presentation hides it. The reader is trained to see a \emph{matrix}: a
covariance in finance, a Gram matrix in regression, a kernel in learning---formed,
stored, and factored.

None of the algorithms in this book sees that matrix. Inspecting them,
the quadratic form enters only through a short list of linear \emph{actions}. So the
quadratic form is, for the algorithm, an \emph{operator}: defined by what it does to a
vector, not by a table of $n^2$ numbers. That change of model is what turns one
algorithm into a method for many problems, because the operator may be supplied in
whatever form is cheapest.

This chapter states the contract, gives the four backends that recur, works out the
cost of the most important one, and says where the crossover falls.

\section{The contract}

\begin{definition}[Quadratic-form operator]
\label{def:qform}
\index{operator!quadratic-form contract}
An implementation of the quadratic form $H$ supplies three linear actions:
\begin{enumerate}
\item \texttt{matvec}: $v \mapsto Hv$;
\item \texttt{solve\_free}: $R \mapsto H_{\Free\Free}^{-1}R$ for a multi-column
      right-hand side $R$ and a given free set $\Free$;
\item \texttt{cross}: $x_\Bound \mapsto H_{\Free\Bound}\,x_\Bound$.
\end{enumerate}
\end{definition}

That is the whole interface, and it is worth checking against the algorithms already
seen. The dual method of Chapter~\ref{ch:walking} carries $J$ with $JJ\T = G^{-1}$ and
consumes it only through products---so its needs are a superset of \texttt{matvec} but
the same in kind. The block-pivoting loop of Chapter~\ref{ch:guessing} needs
\texttt{solve\_free} for the inner system and \texttt{cross} for the reduced gradient
at bound indices; Section~\ref{sec:operator-two} of Chapter~\ref{ch:krylov} named
exactly those two. The homotopy of Chapter~\ref{ch:parametric} needs all three: the
bordered solve is \texttt{solve\_free}, the blocked contribution
$H_{\Free\Bound}x_\Bound$ on the right-hand side is \texttt{cross}, and the reduced
gradients that drive the $\mathrm{D}_1$ events are \texttt{matvec}.

\emph{One interface serves the whole book.} That is not a coincidence; it is what the
active-set structure permits, and it is the technical content of the claim in
Chapter~\ref{ch:homotopy} that one engine traces five literatures' paths.

\section{Four backends}

\subsection{Dense}

An explicit symmetric $n\times n$ array. \texttt{matvec} is a matrix--vector product;
\texttt{cross} is an off-diagonal block product; \texttt{solve\_free} is a Cholesky
factorisation of the principal submatrix $H_{\Free\Free}$, falling back to a general
solve on the rare block that is not numerically positive definite. This is the
reference backend and the one against which the others are validated.

Cost: $O(n^2)$ storage, $O(n_\Free^3)$ per free-block solve.

\subsection{Incremental dense}

The same explicit matrix, but \emph{maintaining} $H_{\Free\Free}^{-1}$ rather than
refactorising at every step. Because an active-set method changes the free set by one
index between consecutive solves, a rank-one bordered update (on entry) or a symmetric
deletion (on exit) modifies the stored inverse in place at $O(n_\Free^2)$, against
$O(n_\Free^3)$ for a fresh solve.

The trade is numerical, and it is worth stating plainly because it is the recurring
shape of this decision. A maintained inverse accumulates floating-point error over the
$O(n)$ updates of a trace, whereas a fresh solve is clean at every step. A guard is
needed---a non-positive or non-finite pivot, or any free-set change that is not a
single-index flip, triggers a full refactorisation---and even with it, on
ill-conditioned or near-degenerate inputs the plain dense backend is the safer choice.
So this is an opt-in fast path, not a default. Roughly a factor of two on a
well-conditioned dense problem of a few hundred variables.

\subsection{Factor: diagonal-plus-low-rank}
\label{sec:factor}

\begin{equation}
\label{eq:factor}
H \;=\; \diag(d) \;+\; U\,\Delta\,U\T,
\qquad d \in \R^{n}_{>0},\; U \in \R^{n\times k},\; \Delta \in \R^{k\times k},
\end{equation}
with $k \ll n$. This is the form taken by a factor risk model \cite{sharpe1963}, by a
random-matrix-cleaned covariance (Chapter~\ref{ch:conditioning}), by an approximate
factor model \cite{chamberlain1983} and its modern estimators
\cite{fan2013poet}, and by a low-rank kernel approximation. One object, arrived at
independently by finance, statistics, and machine learning.
\index{factor model}

The free-block solve never forms $H_{\Free\Free}$. By Proposition~\ref{prop:woodbury},
\begin{equation}
\label{eq:woodbury}
H_{\Free\Free}^{-1} \;=\; D_\Free^{-1} - D_\Free^{-1} U_\Free\, W^{-1}\, U_\Free\T D_\Free^{-1},
\qquad
W \;=\; \Delta^{-1} + U_\Free\T D_\Free^{-1} U_\Free \;\in\;\R^{k\times k},
\end{equation}
so a solve costs $O(n_\Free k^2 + k^3)$ rather than $O(n_\Free^3)$, and no $n\times n$
matrix is ever allocated: memory and per-solve cost scale as $O(nk)$ instead of
$O(n^2)$. The cross product contributes only through the low-rank part, the diagonal
having no off-diagonal block:
$H_{\Free\Bound}x_\Bound = U_\Free\Delta(U_\Bound\T x_\Bound)$.

\subsection{Gram: straight from the data matrix}
\label{sec:gram}

A sample covariance \emph{is} a Gram matrix. If $X_c \in \R^{T\times n}$ is the
column-centred data,
\begin{equation}
\label{eq:gram}
H \;=\; \frac{1}{T-1}\,X_c\T X_c,
\end{equation}
and the contract is implemented straight from $X_c$: \texttt{matvec} is the pair of
thin products $X_c\T(X_c v)/(T-1)$, \texttt{cross} restricts those products to free
and blocked columns, and \texttt{solve\_free} works on $X_{c,\Free}$ alone. Memory is
$O(Tn)$ rather than $O(n^2)$. This is the backend Chapter~\ref{ch:krylov} runs CG
against.

The catch is rank: $X_c\T X_c$ has rank at most $T-1$. When $T \ge n$ with full column
rank, $H \succ 0$ and this is the dense reference at lower memory. In the
short-sample regime $T < n$---a covariance estimated from far fewer periods than
assets---$H$ is singular and the free block becomes unsolvable once the free set
outgrows the data rank. An optional ridge $\delta I$ restores positive definiteness,
and the free-block solve is then \eqref{eq:woodbury} carried out in the
$T$-dimensional observation space, at $O(n_\Free T^2 + T^3)$. \emph{This is precisely a
factor model whose factors are the observations themselves}, with $U = X_c\T$ and
$k = T$.

\begin{remark}[A memory benefit, not an unconditional speed-up]
When $T \ge n$ and the dense matrix fits in memory, the dense backend is faster per
step: it slices a precomputed $H_{\Free\Free}$, whereas the Gram backend re-forms
$X_{c,\Free}\T X_{c,\Free}$ at each solve. The same caveat governs the factor backend:
a structured solve pays only while $k$ (or here $T$) stays well below $n$. At full
rank the structured route is \emph{slower} than the plain dense matrix. State this
clearly in any benchmark you report.
\end{remark}

\subsection{Kernels, and sparse designs}

Two more, briefly, because they cost nothing to add. A kernel Gram matrix is given
implicitly by its kernel function; the support-vector-machine path of
Chapter~\ref{ch:homotopy} touches it only through the same actions, so the kernel trick
enters unchanged and the feature space may be infinite-dimensional without the
algorithm noticing. And when $X$ is sparse, the action $v \mapsto X\T(Xv)$ costs the
number of nonzeros rather than $n^2$; the operator inherits the data's sparsity with no
special handling.

\section{The crossover}

When is the structured solve actually cheaper? The question has a clean answer for the
scalar-identity-plus-low-rank case, which is the one Chapter~\ref{ch:conditioning}
produces.

\begin{proposition}[Woodbury versus assemble-and-factor]
\label{prop:crossover}
\index{crossover!Woodbury versus Cholesky}
Assembling $H_{\Free\Free}$ by a rank-$k$ update and factoring it costs
$n_a^2 k + n_a^3/3$ flops, where $n_a = |\Free|$; forming and factoring the $k\times k$
correction $W$ of~\eqref{eq:woodbury} costs $n_a k^2 + k^3/3$. Writing $x = k/n_a$, the
Woodbury solve is cheaper per step iff
$n_a k^2 + \tfrac{k^3}{3} < n_a^2 k + \tfrac{n_a^3}{3}$, which on dividing by $n_a^3$
factors as $(x-1)(x^2 + 4x + 1) < 0$, i.e.\ iff $x < 1$, that is $k < n_a$.
\end{proposition}

The condition $k < n_a$ is embarrassingly mild, and it holds across essentially the
whole of a long-only frontier. On S\&P~500 data, $k = 17$ signal factors against
$n_a = 46$ assets held at the minimum-variance solution, so $x \approx 0.37$. On a
synthetic frontier $x$ runs from about $0.33$ at the high-return end down to about $1$
at the tightest point, where the number of holdings dips to roughly $k$ itself. Only at
that extreme are the two comparable; elsewhere the flop advantage applies, and the
measured gap is larger still---consistent with the $k\times k$ working set fitting in
lower cache levels than an assembled $n_a \times n_a$ system.

Two practical notes. When $n_a \le k$ a solver should fall back to the assemble branch.
And the Woodbury solve itself remains \emph{valid} for all $n_a \ge 1$, since
$W = \Delta^{-1} + \bar\mu^{-1}U_\Free\T U_\Free$ is a positive definite diagonal plus
a positive semidefinite matrix, hence always invertible; only its cost advantage
disappears.

\begin{remark}[Sharing a factorisation across right-hand sides]
The $p+1$ right-hand sides of an equality-augmented solve
(Chapter~\ref{ch:krylov}, Section~3) all share the single factorisation of $W$, so each
right-hand side beyond the first costs only $O(n_a k)$. Structure the code so that
$W$ is factored once per outer step and the solves are applied to a block, not looped.
This is the difference between $p+1$ solves and one solve with $p+1$ columns, and on a
factor model it is the whole of the cost.
\end{remark}

\section{Regularisation through the operator}
\label{sec:reg-operator}

A regularised objective is supported whenever the penalty is \emph{quadratic}, and for
exactly the reason the bare algorithms work: a quadratic penalty changes only the
Hessian, so the systems stay linear and---in the parametric case---the path stays
piecewise linear. Ridge is the canonical case,
\begin{equation}
\label{eq:ridge}
\tfrac12\, w\T H\, w - \lambda\,\mu\T w + \tfrac{\gamma}{2}\,\norm{w}_2^2
\;=\;
\tfrac12\, w\T(H + \gamma I)\, w - \lambda\,\mu\T w,
\end{equation}
which is the substitution $H \mapsto H + \gamma I$; more generally $H + \gamma M$ for
any $M \succeq 0$---a factor-exposure penalty, a deviation-from-benchmark term, a
quadratic turnover penalty $\norm{w - w_{\mathrm{prev}}}_2^2$.

Because the algorithms reach the quadratic form only through
Definition~\ref{def:qform}, the regularised form is just another operator: the penalty
is supplied by the \emph{backend} and the loop is untouched. The result is the exact
solution, or the exact path, of the regularised objective. Two backends already realise
this: the Gram backend takes a ridge directly, and the factor
form~\eqref{eq:factor} is itself a structured ridge, positive definite by the
idiosyncratic floor $d > 0$.

The boundary follows from the same requirement, and is worth knowing.

\begin{itemize}
\item Shrinkage toward a \emph{nonzero} target $w_0$ is still quadratic but adds a
      $\lambda$-independent linear term $-\gamma w_0\T w$, which an interface whose
      linear term is exactly $-\lambda\mu$ does not expose. Fixable, but not free.
\item A nonsmooth $\ell_1$ penalty is not a term in this objective at all. For the
      long-only fully invested book it is vacuous
      (Proposition~\ref{prop:hinge}); where it bites, it is traced by the LASSO
      homotopy of Chapter~\ref{ch:homotopy}; and a gross-exposure \emph{constraint}
      $\norm{w}_1 \le c$ is polyhedral and expressible as inequality rows.
\item Genuinely non-quadratic penalties---entropy, a log barrier---break the structure
      and lie outside the method.
\end{itemize}

\section{What the operator reading is for}

Two claims, and they are the reason this chapter exists.

\paragraph{Cost.} In every case the algorithm is the same loop and returns the
identical answer; only the cost of a step changes, from $O(n^2)$ to $O(nk)$ or to the
number of nonzeros. That is how a method designed in 1956 for a dense covariance
reaches universe sizes its dense formulation could not.

\paragraph{Transfer.} Seeing the members as one says that the factor risk model and the
kernel trick are the \emph{same device}---an operator standing in for a matrix---so
either field's machinery for it is immediately available to the other. And the traffic
is asymmetric, because each field has matured a tool the others under-use. In finance
the quadratic form is almost always a \emph{modelled} object, a diagonal-plus-low-rank
risk model being the industry default, since forming and trusting a dense covariance at
scale is infeasible; sparse regression, kernels aside, takes the Gram matrix as whatever
the design delivers and seldom models it as structured. The habit of replacing the
quadratic form by a cheap structured operator is one that regression can import
wholesale from finance. The traffic runs the other way too: inequality constraints and
the path's inferential apparatus are routine in statistics and absent from the canonical
portfolio codes.

Handed a factor covariance together with $X\T y$, the same engine traces a least angle
regression path in $O(nk)$ per step; under a linear inequality it traces a constrained
path that standard LARS implementations do not provide. These are the under-occupied
cells of the design space of Section~\ref{sec:taxonomy}: the same engine with a
different operator and one extra event type, not new algorithms.

\begin{notes}
The operator contract of Definition~\ref{def:qform} is the \texttt{QuadraticForm}
protocol of the \texttt{cvxcla} package \cite{cvxcla}, whose four backends are those of
Section~2; the reading of it as a cross-field device is
\cite{schmelzer2026perspective}. Proposition~\ref{prop:crossover} and its empirical
crossover values are from \cite{schmelzer2026rmt}. The Woodbury identity is standard
\cite{golub2013}; the diagonal-plus-low-rank covariance goes back to Sharpe's
single-index model \cite{sharpe1963} and the approximate factor model of Chamberlain
and Rothschild \cite{chamberlain1983}, with POET \cite{fan2013poet} a modern estimator
returning the same form.
\end{notes}

\begin{exercises}

\begin{exercise}
For each algorithm in Chapters~\ref{ch:walking}, \ref{ch:guessing}, \ref{ch:krylov},
and~\ref{ch:parametric}, list exactly which of the three actions of
Definition~\ref{def:qform} it uses and where.
\end{exercise}

\begin{exercise}
Verify~\eqref{eq:woodbury} directly by multiplying out
$H_{\Free\Free}\,H_{\Free\Free}^{-1}$, and check that $W \succ 0$ whenever
$\Delta \succ 0$ and $d > 0$.
\end{exercise}

\begin{exercise}
Show that the Gram backend with a ridge is the factor form~\eqref{eq:factor} with
$U = X_c\T$, $k = T$, $\Delta = I/(T-1)$, $d = \delta\onevec$. What does
Proposition~\ref{prop:crossover} then say about when it pays?
\end{exercise}

\begin{exercise}
Count the flops of the two branches in Proposition~\ref{prop:crossover} for
$n_a = 500$ and $k \in \{5, 50, 200, 500, 800\}$, and confirm the crossover at $k = n_a$.
\end{exercise}

\begin{exercise}
Show that $(H_{\Free\Free})^{-1} \ne (H^{-1})_{\Free\Free}$ in general, and give a
$2\times2$ counterexample. Which of the three actions does this fact complicate, and
why does it matter for preconditioning (Chapter~\ref{ch:conditioning})?
\end{exercise}

\begin{exercise}
Explain why a quadratic turnover penalty $\norm{w - w_{\mathrm{prev}}}_2^2$ is
expressible through the operator but shrinkage to a nonzero target is not, given the
interface's fixed linear term.
\end{exercise}

\begin{exercise}[Implementation]
Implement Definition~\ref{def:qform} as an interface with the dense and factor
backends. Verify that a solver written against the interface returns identical answers
from both when the factor backend is handed $\diag(d) + U\Delta U\T$ and the dense
backend is handed the same matrix formed explicitly. Report the agreement.
\end{exercise}

\begin{exercise}[Implementation]
Time both backends over a path trace for $n \in \{200, 500, 1000, 2000\}$ with $k = 20$
fixed, and plot per-step cost against $n$. Then hold $n = 500$ and vary $k$ to locate
the crossover empirically. Does it fall where Proposition~\ref{prop:crossover} predicts,
and if not, why not?
\end{exercise}

\end{exercises}
